% sec_intro.tex — Introduction
\section{Introduction}\label{sec:intro}

Autonomous systems are increasingly deployed to act persistently on behalf of
their principals: LLM-based agents that plan, remember, and use tools
\citep{wang2023,xi2023,park2023}; recommender and notification systems that
continuously shape what users see; rehabilitation and coaching devices that
nudge human habits. In each deployment the same two questions arise, and both
are currently unanswerable in operational terms: \emph{how much autonomy does
the agent retain under a given influence}, and \emph{how much did a given
interaction take away}. Governance frameworks grade agents by levels and
profiles \citep{feng2025,kasirzadeh2025,dung2024}; psychology measures
motivation and control \citep{ryan_deci2000,maier_seligman2016}; information
theory measures control loops \citep{klyubin2005,massey1990}; AI safety
catalogues failure modes one by one
\citep{amodei2016,everitt2021,greshake2023}. What is missing is a common
yardstick: a measure of autonomy itself, defined independently of substrate,
decomposed into auditable channels, and governed by explicit axioms and
theorems rather than by unarticulated conventions.

Four misconceptions stand in the way of such a measure, and we dispel them at
the outset. Autonomy is \emph{not isolation}: a closed system is dead, not
free; adaptation requires informational intake (our Axiom~AX2). Autonomy is
\emph{not randomness}: to be autonomous is to be determined by one's self, not
by nothing (Axiom~AX0 makes directed information from internal states to
actions the defining positive quantity). Autonomy is \emph{not omnipotence}:
capability is one channel among five, and its absence does not annihilate the
others. And autonomy is \emph{not the absence of relations}: influences divide
into constitutive support (our $\sigma$ class) and reductive interference;
``zero relations'' is not the maximum of autonomy but lies outside its feasible
region \citep{mackenzie_stoljar2000}.

\paragraph{The proposal.} We define the autonomy of an agent relative to an
environment and a time horizon as the degree to which its behavior is
determined by its own normative structure $(G,B,M)$---its goals, beliefs, and
self-model---in five respects: \emph{origination} (behavior is causally sourced
in internal states rather than bypassed), \emph{informedness} (external
influences are registered by the self-model), \emph{endorsability} (changes to
the motivational structure pass, or could pass, second-order review that is
itself free of tampering---a historical clause), \emph{executability} (intent
has predictable leverage on the situation, with exit options), and
\emph{viability} (the resource reproduction sustaining all of the above is not
interrupted). Each clause becomes a channel with an estimator, a canonical
ratio scale, and a zero semantics; a gating viability factor multiplies a
composite of the five. The resulting scale is comparable \emph{within a layer
at a fixed instantiation}; cross-layer and cross-substrate readings are
reported as channel profiles (Section~\ref{sec:framework}).

\paragraph{A note on objectivity, at the outset.} Any measure of autonomy is
built from indicators, and every indicator carries its designer's limitations;
no instrument family is neutral. We therefore do not claim that our channels
\emph{are} autonomy, nor that subjectivity can be designed away. What can be
built is \emph{procedural objectivity}: every design commitment is registered
and priced (including what is lost if it is rejected); invariance under
irrelevant re-instrumentation is itself made an axiom; the boundary between
the knowable and the unknowable is proved rather than assumed; and the
channel choice itself is exposed to a registered external falsifier. The
measure is objective in the only sense measurement science delivers---not
free of conventions, but explicit about them and stable under independent
reconstruction (Section~\ref{sec:objectivity}).

\paragraph{Contributions.}
The primary contribution is the \emph{audit schema} itself---a measurement
and identifiability discipline for autonomy under influence; the theorems
are its load-bearing instruments, not a claim to a finished mathematical
theory of autonomy.
\begin{enumerate}
\item \textbf{A formal framework} (Section~\ref{sec:framework}): a
time-unfolded causal graph with a five-class intervention taxonomy
($\sigma/\alpha/\beta/\gamma/\delta$) whose completeness is a theorem
relative to that interface (Proposition~\ref{thm:taxonomy-completeness}):
within the graph's surgery loci, every external influence acts by changing
what the agent sees, its mechanism itself, what it can do, or its conditions
of self-maintenance. The $\alpha/\gamma$ boundary is a structural criterion
(the intact channel survives as a component, or not), and the one deliberate
boundary---surgeries on the environment's own dynamics---is registered as
re-instantiation of the measurement, not a sixth class.
\item \textbf{Measures from axioms} (Sections~\ref{sec:axioms} and
\ref{sec:representation}): six channel measures, each with an absolute zero and
a canonical normalizer; and a representation theorem (Theorem~\ref{thm:rep})
showing that continuity, strict coordinatewise monotonicity, homogeneity, and
coordinatewise scale invariance of the induced order admit exactly the weighted
geometric mean---``forced'' relative to axioms that are defended normative
commitments, not measurement-theoretic necessities (the scale-invariance
characterization itself is classical: \citealp{ebert_welsch2004}; our additions
are scoped in Section~\ref{sec:representation})---with two derivations,
aggregation-theoretic and measurement-theoretic, independent in toolchain
though sharing the invariance axioms, and a harmonic-mean counterexample
(Proposition~\ref{prop:harmonic}) showing exactly why the invariance axiom is
indispensable. Biseparability and zero annihilation are consequences of the
minimal invariance trio, not independent commitments: the two derivations
pin the minimal sufficient set at $\{$AX4, AX6b, AX8$\}$. The load-bearing
axioms themselves are normative commitments of the framework, defended in
Section~\ref{sec:objections} rather than derived.
\item \textbf{A suite of theorems with falsification windows}
(Section~\ref{sec:theorems}): coverage and multi-channel necessity (T2---no
single channel detects all intervention classes; within this instrument
family the decomposition is an attribution cover, with necessity as a
theorem); manipulation strictly reduces autonomy wherever the affected
instrument enters the composite, and covert manipulation
inflates a measurable
hallucination index (T3);
highest-leverage attack concentration and water-filling defense under
log-linear damage (T4);
loss of autonomy as a six-way disjunction (T5); an explicit autonomy margin
$\deltastar=\min(\theta/W,\kappa,\lceil(r{+}1)/2\rceil c_0/W)$ with a
positive retention floor
$\Amin$---including the proof that \emph{detection alone can provide no margin
against slow drift}, so that margins come from repair capacity while
detection's role is bounding the registration domain, and a sharp
blind-zone prediction (T6); component-level delegation laws (T7); and the
coercion paradox as a consistency result (T8).
\item \textbf{Identifiability theory} (Section~\ref{sec:identifiability}):
white-box point identification, tiered by the modal commitments each channel
carries; a certified \emph{one-sided} black-box bound for $\Aorg$ under
bypass-semantics probes (the reference lower end is uncertified) and
protocol-relative bounds for $\Aexec$; and Proposition~\ref{prop:unidentifiable}
(implanted endorsement is behaviorally indistinguishable from endogenous
endorsement), which makes the impossibility of black-box historical auditing a
theorem and full causal logging an architectural necessity.
\item \textbf{A simulation benchmark with registered falsification windows}
(Section~\ref{sec:experiments}): a 4{,}620-run campaign over fifteen
scenarios---six attack classes, delegation, helplessness induction, resource
starvation, capability-envelope erosion, forced autonomy, black-box
degradation, cross-framework consistency, and an execution dynamic-range
arm, plus an adaptive-adversary supplement. The campaign is an
implementation-level consistency check, and we report it as such: all
fifteen proposition-level criteria are met within the registered windows with
zero refutations---the expected outcome of a correct implementation, not
independent evidence for the theory---and the one adverse finding is a
sensitivity calibration of the black-box bound's premise under adaptive
adversaries (BB-2). Verdicts are
stable under alternative aggregation weights (Kendall $\tau\ge0.93$) and
under axiom-relaxation robustness checks in the elasticity parameter; on the
default instance the execution channel's narrow band makes that statement an
effectively four-channel one, so the same weight battery was rerun on the
rich-coupling EXEC-1 instance (fifth channel at full dynamic range): no
verdict flips there either.
\item \textbf{An interpretive correspondence with existing frameworks}
(Proposition~C0,
\S\ref{sec:theorems} and \S\ref{sec:discussion}): under explicit definitional
readings---whose faithfulness to the source texts is registered as
hermeneutic, not mathematical, residual uncertainty---the agency profiles of
\citet{dung2024}, the four dimensions of \citet{kasirzadeh2025}, and the
user roles of \citet{feng2025} are exhibited as projections of our channel
space onto the special case of a benign principal, with a constructive
witness that our state descriptor separates configurations they cannot
(Proposition~C0(ii)). This is an explanatory correspondence under registered
readings, not a mathematical unification theorem.
\end{enumerate}

\paragraph{Map of the paper.} Section~\ref{sec:related-work} positions the
framework in seven literatures and states the gaps it fills.
Section~\ref{sec:framework} gives the operational definition, the system
model, and the intervention taxonomy. Section~\ref{sec:axioms} states the axioms and
the six channel measures. Section~\ref{sec:representation} proves the
representation theorem. Section~\ref{sec:theorems} presents the theorem suite
with proof sketches and falsification windows; full proofs are in the
appendices. Section~\ref{sec:identifiability} develops the identifiability
analysis. Section~\ref{sec:experiments} reports the simulation study.
Section~\ref{sec:discussion} answers the standing debates and registers
limitations; Section~\ref{sec:conclusion} concludes.
